% Speculative decoding: a cheap draft proposes tokens serially, while one
% target-model pass verifies their conditional probabilities in parallel.
\documentclass[border=2mm]{standalone}
\usepackage{amsmath,amssymb}
% Chinese labels in figures; fonts are shared with the book.
\usepackage[UTF8,fontset=none]{ctex}
% CJK fonts for the book, the figures and the slides.
% Noto Sans CJK SC (SIL Open Font License) is the same typeface as Source Han Sans SC.
% macOS: brew install --cask font-noto-sans-cjk-sc
% Debian/Ubuntu: apt install fonts-noto-cjk
\setCJKmainfont[BoldFont={Noto Sans CJK SC Bold}]{Noto Sans CJK SC}
\setCJKsansfont[BoldFont={Noto Sans CJK SC Bold}]{Noto Sans CJK SC}
\setCJKmonofont{Noto Sans CJK SC}


% Shared visual language for the figures of the book.
% Every figure in this directory loads this file so the chapter reads as one system.
% NOTE: no blank lines are allowed inside \tikzset{...} (a blank line emits \par
% and breaks key scanning).
\NeedsTeXFormat{LaTeX2e}
\ProvidesPackage{figstyle}[2026/10/21 Figure style for Free Energy to LLMs]

\RequirePackage{tikz}
\usetikzlibrary{arrows.meta,positioning,calc,fit,shapes.geometric,shapes.misc,%
  backgrounds,matrix,decorations.pathreplacing,decorations.markings,patterns.meta}

% ---------------------------------------------------------------- palette
% Same values as book/preamble/style.tex, prefixed with F.
\definecolor{FNavy}{HTML}{14324A}
\definecolor{FTeal}{HTML}{0E7490}
\definecolor{FSky}{HTML}{E7F5F8}
\definecolor{FOrange}{HTML}{D97706}
\definecolor{FAmber}{HTML}{FFF3D6}
\definecolor{FInk}{HTML}{1F2937}
\definecolor{FSlate}{HTML}{64748B}
\definecolor{FMist}{HTML}{F5F7FA}
\definecolor{FRose}{HTML}{9F1239}
\definecolor{FRoseMist}{HTML}{FFF1F2}

% ---------------------------------------------------------------- semantics
% Teal   = probability / model / data object
% Orange = training objective, cost, or the quantity moved through the system
% Navy   = neutral structure
% Slate  = secondary label or metadata
% Rose   = failure mode, hazard, masked-out region
\tikzset{
  % Lay figures out in millimetres.  A figure designed 168mm wide is included at
  % \textwidth, so its fonts appear at exactly the size written here.
  every picture/.style={x=1mm,y=1mm,font=\small,line cap=round,line join=round},
  % ---- block styles
  nbase/.style={draw,rounded corners=1.2pt,align=center,inner sep=2.4pt,%
    minimum height=7mm,text=FInk},
  nteal/.style={nbase,draw=FTeal,fill=FSky},
  norange/.style={nbase,draw=FOrange,fill=FAmber},
  nnavy/.style={nbase,draw=FNavy,fill=white},
  nsoft/.style={nbase,draw=FSlate!55,fill=FMist},
  nrose/.style={nbase,draw=FRose,fill=FRoseMist},
  ghost/.style={draw=none,fill=none,align=center,inner sep=1pt},
  plain/.style={draw=none,fill=none,align=center,inner sep=0pt},
  % ---- arrows
  ar/.style={-{Latex[length=1.7mm,width=1.35mm]},draw=FSlate,line width=0.75pt},
  ard/.style={ar,dashed},
  arteal/.style={ar,draw=FTeal},
  arorange/.style={ar,draw=FOrange},
  arrose/.style={ar,draw=FRose},
  thickar/.style={ar,line width=1.2pt},
  % ---- text roles
  ftitle/.style={font=\large\bfseries,text=FNavy,align=center},
  ptitle/.style={font=\small\bfseries,text=FNavy,align=center},
  ptitleo/.style={font=\small\bfseries,text=FOrange,align=center},
  lbl/.style={font=\scriptsize,text=FSlate,align=center,inner sep=1pt},
  lblt/.style={lbl,text=FTeal},
  lblo/.style={lbl,text=FOrange},
  lbln/.style={lbl,text=FNavy},
  lblr/.style={lbl,text=FRose},
  sub/.style={font=\scriptsize,text=FSlate!90,align=center,inner sep=1pt},
  monob/.style={font=\ttfamily\scriptsize,text=FInk,inner sep=1pt},
  mathb/.style={font=\small,text=FInk,inner sep=1pt},
  eqn/.style={font=\scriptsize,inner sep=1.5pt,text=FInk},
  % ---- decorations
  brace/.style={decorate,decoration={brace,amplitude=2.4pt},draw=FSlate,%
    line width=0.6pt},
  bracet/.style={brace,draw=FTeal},
  braceo/.style={brace,draw=FOrange},
  axis/.style={draw=FSlate!55,line width=0.6pt},
  guide/.style={draw=FSlate!45,line width=0.5pt,dashed},
  halolabel/.style={font=\scriptsize,text=FInk,fill=white,fill opacity=0.85,%
    text opacity=1,inner sep=1.4pt,align=center},
}

% Figure title above the drawing, anchored at the figure's top-left.
\newcommand{\figtitle}[2][0]{%
  \node[ftitle,anchor=north west,outer sep=0pt] at (#1,0) {#2};}

% Rounded panel behind a subgroup of nodes.
\newcommand{\figpanel}[3][]{%
  \node[draw=FSlate!35,fill=#3,rounded corners=2pt,inner sep=2.6mm,#1] {#2};}

\endinput

\begin{document}
\begin{tikzpicture}
  \node[ptitle,anchor=west] at (0,78) {推测解码：草稿模型提议，目标模型一次验证};

  % Draft proposals are generated one after another.
  \node[lbln] at (82,72) {逐个生成候选};
  \node[nteal,minimum width=25mm,minimum height=9mm] (draft) at (15,64) {草稿模型 $q$};
  \foreach \x/\i in {48/1,68/2,88/3,108/4} {
    \node[norange,minimum width=15mm,minimum height=9mm] (x\i) at (\x,64) {$x_\i$};
  }
  \draw[arorange] (draft.east)--(x1.west);
  \draw[arorange] (x1.east)--(x2.west);
  \draw[arorange] (x2.east)--(x3.west);
  \draw[arorange] (x3.east)--(x4.west);

  % A single causal target pass evaluates all proposed positions.
  \node[nteal,minimum width=81mm,minimum height=12mm] (target) at (78,45)
    {目标模型 $p_\theta$\\一次前向验证四个位置};
  \foreach \x in {48,68,88,108} {
    \draw[arteal] (\x,59.4)--(\x,51.6);
    \draw[ar] (\x,38.8)--(\x,33.6);
  }

  % The acceptance decision is sequential: after the first rejection, later
  % draft tokens are discarded even though the target pass evaluated them.
  \node[text=FTeal,font=\Large] at (48,29.5) {$\checkmark$};
  \node[text=FTeal,font=\Large] at (68,29.5) {$\checkmark$};
  \node[text=FRose,font=\Large] at (88,29.5) {$\times$};
  \node[text=FSlate,font=\Large] at (108,29.5) {$\times$};
  \node[lblt] at (48,22) {接受};
  \node[lblt] at (68,22) {接受};
  \node[lblr] at (88,22) {首次拒绝};
  \node[lbl] at (108,22) {后续丢弃};
  \draw[bracet] (40,17)--(76,17);
  \node[lblt] at (58,12.5) {保留前缀 $x_1,x_2$};
  \draw[arrose] (88,17)--(88,13);
  \node[nrose,minimum width=24mm,minimum height=8mm] at (88,7.5) {修正采样};
  \node[lbln] at (19,7.5) {目标模型仅做一次前向};

  % Comparison: ordinary decoding needs a separate target pass per token.
  \draw[FSlate!40,line width=.6pt] (125,5)--(125,74);
  \node[lbln] at (147,72) {普通逐 token 解码};
  \foreach \y/\i in {61/1,47/2,33/3,19/4} {
    \node[nsoft,minimum width=31mm,minimum height=8mm] (p\i) at (147,\y)
      {$p_\theta\ \to\ x_\i$};
  }
  \draw[ar] (p1.south)--(p2.north);
  \draw[ar] (p2.south)--(p3.north);
  \draw[ar] (p3.south)--(p4.north);
  \node[lbl] at (147,8) {四次串行前向};
\end{tikzpicture}
\end{document}
